\documentclass[12pt]{article}
\usepackage[utf8]{inputenc}
\usepackage[T1]{fontenc}
\usepackage[margin=1in]{geometry}
\usepackage{setspace}
\usepackage{xcolor}
\usepackage[colorlinks=true, urlcolor=blue!60!black, linkcolor=blue!60!black]{hyperref}
\setlength{\parindent}{2em}
\setlength{\parskip}{1.25em}
\renewcommand{\baselinestretch}{1.0}
\begin{document}
\begin{center}
{\LARGE \textit{Trump v. Mazars USA, LLP}, 591 U.S. 848 (2020)\par}
\vspace{0.8em}
{\large Case Brief --- Four-Page Edition\par}
\end{center}
\vspace{1em}

\section{Memory Jogger}

When House committees subpoenaed a sitting President's personal financial records from his accountants and banks, the Court held 7--2 that neither the President's heightened executive-privilege standard nor the House's deferential standard governs; courts must weigh four separation-of-powers considerations before enforcing such a subpoena (\href{https://supreme.justia.com/cases/federal/us/591/19-715/}{Justia}).

\section{Detailed Case Facts}

In April 2019 three House committees issued four subpoenas for roughly a decade of financial records concerning President Trump, his children, and affiliated businesses. The Financial Services Committee subpoenaed Deutsche Bank and Capital One, citing anti-money-laundering, counterterrorism, and real-estate legislation. The Intelligence Committee subpoenaed Deutsche Bank in an inquiry into Russian interference and foreign financial leverage. The Oversight Committee subpoenaed the President's accounting firm, Mazars USA, for 2011--2018 financial statements and related documents after Michael Cohen's testimony, citing possible illegal conduct, undisclosed conflicts, the Emoluments Clauses, and the accuracy of financial disclosures.

\section{Procedural History}

The President, suing personally, and his family and businesses sued in the District of Columbia (\textit{Mazars}) and the Southern District of New York (\textit{Deutsche Bank}). They did not assert executive privilege but argued that the subpoenas lacked a legitimate legislative purpose. The D.C. Circuit affirmed judgment for the House over Judge Rao's dissent, which would channel investigation of presidential wrongdoing to impeachment. The Second Circuit largely affirmed, reasoning that the records were personal and the President could serve as a ``case study,'' over Judge Livingston's partial dissent. The Supreme Court consolidated the cases.

\section{Judicial Votes}

Vacated and remanded, \textbf{7--2}, on July 9, 2020.
\begin{itemize}
\item \textbf{Majority:} Chief Justice Roberts, joined by Justices Ginsburg, Breyer, Sotomayor, Kagan, Gorsuch, and Kavanaugh.
\item \textbf{Dissenting:} Justice Thomas; Justice Alito.
\end{itemize}

\section{Holding}

The lower courts failed to account for the separation-of-powers concerns raised by congressional subpoenas for a sitting President's personal information. Rejecting both the heightened standard the President urged and the House's deferential legislative-purpose standard, the Court required a balanced, non-exhaustive four-factor inquiry before such a subpoena is enforced.

\section{Analysis of Opinions}

\textbf{Majority (Roberts).} For two centuries, from Washington's St.\ Clair inquiry through the Clinton administration, the branches resolved document disputes through the ``hurly-burly, the give-and-take of the political process''; under \textit{Noel Canning}, that practice deserves great weight. Congress's implied power ``to secure needed information'' (\textit{McGrain}; \textit{Watkins}) requires a ``valid legislative purpose'' and may not be used for ``law enforcement'' or ``to expose for the sake of exposure.''

The Court refused to ``transplant that protection root and branch'' from \textit{United States v.\ Nixon}'s privileged communications to private records, which would impede legislating. But the House's approach ignored that Congress and the President are ``opposite and rival'' branches: nearly any personal paper could relate to conceivable legislation. Concerns do not vanish because papers are personal or held by third parties, since the Constitution ``deals with substance, not shadows.'' Courts must weigh: (1) whether other sources could supply the information, so the President is not used as a ``case study''; (2) whether the subpoena is no broader than reasonably necessary; (3) the detail and substance of the evidence of legislative purpose; and (4) the burdens on the President.

\textbf{Thomas (dissenting).} Under original meaning, Congress has no implied power to subpoena private, nonofficial documents of anyone; \textit{McGrain} is unmoored from text, and the House must proceed under impeachment.

\textbf{Alito (dissenting).} Such subpoenas are ``inherently suspicious,'' and there was ``disturbing evidence of an improper law enforcement purpose.'' He would require the House to describe the legislation, its authority, the scope, and the lack of alternatives. On remand, the D.C. Circuit (2022) applied the factors after Trump left office and narrowed the Oversight subpoena.

\section{Examples: Future Applications}

\begin{enumerate}
\item \textbf{Same-Side: Family Resort ``Case Study.''} A committee drafting a cap on foreign-government resort payments subpoenas eight years of guest lists for the President's family resort. Other data could suffice, the demand is overbroad, and the purpose evidence is thin.
\item \textbf{Same-Side: Presidential Medical Records.} A committee drafting disability procedures subpoenas ten years of the President's medical records. Expert testimony could inform a general statute, so enforcement should be refused.
\item \textbf{Opposite-Side: State Grand Jury.} A state grand jury subpoenas the President's accountants. Under \textit{Trump v.\ Vance}, decided the same day, criminal process faces no heightened need standard.
\item \textbf{Opposite-Side: Private Executive.} A committee subpoenas a bank chief executive's accounts for bank-capital legislation. Ordinary \textit{McGrain}--\textit{Watkins} review applies.
\item \textbf{Fence-Sitter: Former President.} A narrow subpoena cites specific discrepancies in a former President's disclosure forms. It satisfies several factors, but whether they apply with full force after office is undecided by the Supreme Court.
\end{enumerate}

\section{Critique}

\textbf{Indeterminacy and asymmetry.} Critics note that the non-exhaustive four-factor test offers little bright-line guidance and invites litigation that routinely outlasts a two-year House subpoena, so delay becomes a functional victory. Students should also ask why a county prosecutor in \textit{Vance} faced a lower bar for the same Mazars records than three House committees.

\textbf{The impeachment-channeling problem.} Thomas's dissent raises a structural question: if Congress can compel private presidential papers only through impeachment, it may be pushed toward premature, hyper-partisan impeachment inquiries whenever it needs facts for ethics or disclosure statutes.

\section{Key Quotations}

\begin{quote}
``The President is the only person who alone composes a branch of government. As a result, there is not always a clear line between his personal and official affairs.''
\end{quote}
\noindent --- Chief Justice Roberts, Opinion of the Court.

\begin{quote}
``The President's unique constitutional position means that Congress may not look to him as a `case study' for general legislation.''
\end{quote}
\noindent --- Chief Justice Roberts, Opinion of the Court.

\begin{quote}
``I would hold that Congress has no power to issue a legislative subpoena for private, nonofficial documents---whether they belong to the President or not. Congress may be able to obtain these documents as part of an investigation of the President, but to do so, it must proceed under the impeachment power.''
\end{quote}
\noindent --- Justice Thomas, dissenting.

\end{document}
